The numerator is always positive, so a solution must have $x>-1$. Write $R(x)=x-4+\dfrac{8}{x+1}$; it is positive and has minimum $2\sqrt8-5>\dfrac{1}{2}$. Let $F(x)=\log R(x)-(x-5)\log10$. Since $R^{\prime}(x)=1-\dfrac{8}{(x+1)^2}<1$, if $R^{\prime}\le0$ then $F^{\prime}<0$, and if $R^{\prime}>0$ then $F^{\prime}=\dfrac{R^{\prime}}{R}-\log10<2-\log10<0$. Thus $F$ is strictly decreasing. At $x=5$, $R=\dfrac{7}{3}>1=10^{x-5}$, whereas at $x=6$, $R=\dfrac{22}{7}<10$. Continuity gives a root in $(5,6)$, and strict decrease gives uniqueness. There is $\boxed1$ real solution.
